\section{Guessing the active set instead of walking to it}
\label{sec:pdas}

The dual method reaches the active set one constraint at a time. That is its
defining feature and, at moderate $n$, its principal cost: a budget-plus-bounds
problem in $n = \qpPdasRefN{}$ variables takes some \qpPdasBudgetOuter{} outer
iterations (Section~\ref{ssec:pdas}), each of which is a handful of
matrix--vector products. The alternative is to guess the whole active
set at once, solve the working-set subproblem it induces, and repair the guess
from the signs that come back. This is the semismooth Newton method of
Hinterm\"uller, Ito and Kunisch~\cite{hintermuller2002} read as a primal--dual
active-set strategy, and block principal pivoting on the associated linear
complementarity problem~\cite{judice1994,murty1988,cottle1992} read as a rule for
exchanging several indices at a time.

\subsection{The iteration}

Given a candidate set $\Active$, Proposition~\ref{prop:sub} gives the working-set
solution directly:
\begin{equation}
  \bigl(C_\Active\T G^{-1} C_\Active\bigr)\lambda_\Active
   = b_\Active - C_\Active\T x_u,
  \qquad
  x = x_u + G^{-1}C_\Active \lambda_\Active,
  \qquad x_u = G^{-1}a,
  \label{eq:pdassolve}
\end{equation}
with $\lambda_i = 0$ for $i \notin \Active$. Both solves go through one Cholesky
factorisation of $G$, computed once, and one of the $k \times k$ dual Hessian,
computed per repair. The repair rule reads the two sign conditions of
\eqref{eq:dfeas}--\eqref{eq:pfeas} as instructions:
\begin{equation}
  \Active' = \{1,\dots,m_{\mathrm{eq}}\} \;\cup\;
  \bigl\{\, i > m_{\mathrm{eq}} \;:\;
    (i \in \Active \text{ and } \lambda_i \ge 0)
    \ \text{ or }\
    (c_i\T x < b_i) \,\bigr\},
  \label{eq:repair}
\end{equation}
that is: an active constraint whose multiplier has gone negative does not belong
in the set and leaves it; an inactive constraint that is violated at $x$ does
belong and joins it; equalities are always in. The starting guess is the
equalities together with whatever the unconstrained minimiser violates, which is
already the answer when it violates nothing. If $\Active' = \Active$ the
conditions~\eqref{eq:pfeas}--\eqref{eq:comp} hold to the tolerance used, and
\eqref{eq:pdassolve} supplies \eqref{eq:stat} by construction.

The whole set is exchanged at once, and it converges in a handful of
repairs almost independently of $n$, from \qpPdasRepairsLo{} to
\qpPdasRepairsHi{} across every family and size measured in
Section~\ref{ssec:pdas}. It trades arithmetic, which is abundant at these sizes,
for iterations, which are not.

\subsection{Why there is no finite-termination theorem here}

Block exchange can cycle, and the standard remedy is a least-index rule that flips
only the lowest-indexed offending constraint, in the manner of Bland's rule for
the simplex method and Murty's least-index rule for linear complementarity
problems~\cite{murty1988}. What follows concerns this \emph{block} exchange
specifically, and not active-set methods in general: the exact walk of
Section~\ref{sec:iteration} is finite under nondegeneracy
(Proposition~\ref{prop:ascent}), and for parametric families of QPs its worst-case
iteration count can be certified exactly~\cite{arnstrom2022certification}. In the bound-constrained case that remedy comes with a
theorem. The theorem does not survive the generalisation, and the reason is
structural, not a gap in the argument.

Consider $\min_{x \ge 0} \tfrac12 x\T G x - a\T x$, so that $C = I$ and $m = n$.
The system solved on a candidate set $\Active$ is then
\begin{equation}
  \bigl(I_{\cdot\Active}\T G^{-1} I_{\cdot\Active}\bigr)\lambda_\Active
  = \bigl(G^{-1}\bigr)_{\Active\Active}\lambda_\Active
  = b_\Active - x_{u,\Active},
  \label{eq:principal}
\end{equation}
whose coefficient matrix is a \emph{principal submatrix} of $G^{-1}$. Since
$G^{-1} \succ 0$, every principal submatrix of it is positive definite, so
$G^{-1}$ is a $P$-matrix: every principal pivot is well defined, whatever the
guess. That is the hypothesis under which block principal pivoting with a
least-index guard terminates finitely at the unique
solution~\cite{judice1994,murty1988}, and the guarantee survives inexact inner
solves once their residual is small against the problem's decision
margin~\cite{schmelzer2026nncg}.

None of it is available for general $C$. The matrix
$C_\Active\T G^{-1} C_\Active$ is positive definite exactly when $C_\Active$ has
full column rank, and that is a property of the guess, not of the data: a guessed
set may name $k > n$ constraints, or $k \le n$ linearly dependent ones, and then
the pivot is undefined. There is no $P$-matrix to
appeal to, and correspondingly no theorem to inherit. Contrast
Proposition~\ref{prop:rank}: the dual method's step rules make dependence
unreachable, and guessing forfeits that protection.

The implementation treats the Cholesky factorisation of
$C_\Active\T G^{-1}C_\Active$ as the rank test, where a dependent guess fails
instead of returning something plausible, and the least-index rule as a way of
making progress where block exchange oscillates, and not as a termination proof.
Neither substitutes for the guarantee. The certificate does.

\subsection{The same program in other coordinates}

There is a symmetry here worth making explicit, because it settles what kind of
property the $P$-matrix argument is and, as Section~\ref{ssec:transform} uses it,
supplies problems whose answer is known at any conditioning. It is a different
invariance from Proposition~\ref{prop:invariance}, which fixes what the
factorisation is free to choose; this one fixes what the coordinates are.

\begin{proposition}[Reparametrisation invariance]\label{prop:reparam}
Let $T \in \R^{n \times n}$ be invertible and substitute $x = Ty$, which carries
the data to
\begin{equation}
  G' = T\T G T, \qquad a' = T\T a, \qquad C' = T\T C, \qquad b' = b .
  \label{eq:reparam}
\end{equation}
Then for every index set $\Active$,
\begin{equation}
  C'_\Active{}\T G'^{-1} C'_\Active = C_\Active\T G^{-1} C_\Active,
  \qquad
  C'_\Active{}\T x'_u = C_\Active\T x_u,
  \qquad
  \operatorname{rank} C'_\Active = \operatorname{rank} C_\Active ,
  \label{eq:invdual}
\end{equation}
where $x'_u = G'^{-1}a'$. Consequently the transformed program has the same
feasible slacks at corresponding points, the same optimal active set and the same
multipliers, its solution is $y^\star = T^{-1}x^\star$, and both the walk of
Algorithm~\ref{alg:gi} and the repair rule~\eqref{eq:repair} visit the same sets
in the same order.
\end{proposition}

\begin{proof}
$G'^{-1} = T^{-1}G^{-1}T^{-\top}$, so
$C'_\Active{}\T G'^{-1}C'_\Active
 = C_\Active\T T\,T^{-1} G^{-1} T^{-\top}\,T\T C_\Active
 = C_\Active\T G^{-1} C_\Active$, and $x'_u = T^{-1}G^{-1}T^{-\top}T\T a
= T^{-1}x_u$ gives $C'_\Active{}\T x'_u = C_\Active\T T\,T^{-1}x_u
= C_\Active\T x_u$. The rank claim is immediate from $T\T$ being invertible.
Slacks are invariant, since $C'{}\T y - b = C\T Ty - b = C\T x - b$; so
therefore is every sign test in \eqref{eq:repair}. For the multipliers,
$G'y - a' = C'\lambda$ reads $T\T(Gx - a - C\lambda) = 0$, which holds exactly
when $Gx - a = C\lambda$. Both the coefficient matrix and the right-hand side
of~\eqref{eq:pdassolve} being invariant by \eqref{eq:invdual}, so is
$\lambda_\Active$, and $x = x_u + G^{-1}C_\Active\lambda_\Active$ maps to
$T^{-1}x$ by the same identity. The quantities the iterations branch on are thus
the same in both coordinate systems, and so are the sequences of sets.
\end{proof}

Two things follow, in opposite directions.

\begin{remark}[Rank deficiency is geometric, not a matter of coordinates]
Because $\operatorname{rank}C_\Active$ is preserved, no change of variables can
create a dependent guess where there was none, and none can remove one. The
exposure of this section is therefore a property of the constraint geometry and
not of how $C$ happens to be written. That cuts both ways. A dense $C'$ carrying
no visible structure may be a set of bounds in disguise --- $C' = T\T$ for some
$T$, in which case $C'_\Active{}\T G'^{-1}C'_\Active$ is a principal submatrix of
$G^{-1}$ for the untransformed $G$ and the $P$-matrix argument applies in full,
which nothing in the entries of $C'$ reveals. And the duplicated columns of
Section~\ref{ssec:pdas}'s fifth family cannot be transformed away.
\end{remark}

\begin{remark}[There is nothing here to precondition]
Proposition~\ref{prop:reparam} holds for every invertible $T$, in particular
for the whitening $T = L^{-\top}$ with $G = LL\T$, which makes $G' = I$. So the
iteration count cannot be improved by a change of variables, and the conditioning
of $G$ --- which a Krylov method's iteration count answers to directly --- does
not enter this one's at all. That is not a coincidence but a restatement of
\eqref{eq:inv1}: the factor $J$ with $J\T G J = I$ \emph{is} the whitening
transformation, computed once and carried by the iteration, so a caller who
applies one first has the solver do the same work twice.
Section~\ref{ssec:transform} measures both halves of this, and the arithmetic
loses either way, since the only thing a transformation can change is the
structure of the data and dense is the worst it can be.
\end{remark}

\subsection{The certificate}

Because the method may fail, and may fail by stabilising on a set that is simply not
optimal, the trade is sound only because the answer is checked. Here
Proposition~\ref{prop:sufficient} carries the argument: a candidate
$(x,\lambda)$ satisfying
\begin{equation}
  \norm{G x - a - C\lambda}_\infty \le \varepsilon,
  \quad |c_i\T x - b_i| \le \varepsilon \ (i \le m_{\mathrm{eq}}),
  \quad c_i\T x - b_i \ge -\varepsilon, \ \ \lambda_i \ge -\varepsilon,
  \ \ |\lambda_i (c_i\T x - b_i)| \le \varepsilon
  \label{eq:certificate}
\end{equation}
on the inequalities, with $\varepsilon$ scaled by the data, \emph{is} the answer to
that tolerance.
Stationarity is checked against $G$ directly and not trusted from the
construction, since the construction is what an ill-conditioned working set
corrupts. A candidate that fails is discarded and the exact walk of
Algorithm~\ref{alg:gi} runs instead, so guessing can never degrade the answer: it
produces the minimiser or nothing.

\begin{remark}[Two tolerances with very different jobs]
The tolerance in~\eqref{eq:repair} decides which set is tried next. A bad choice
there costs a repair or a fallback and can never cost correctness, so it is not
delicate. The tolerance in~\eqref{eq:certificate} is the only thing standing
between a non-optimal point and the caller, and is therefore the one to reason
about. It is nonetheless not delicate either, because the quantity it thresholds
is bimodal: a candidate constructed from a well-conditioned working set satisfies
\eqref{eq:certificate} to a few multiples of the machine epsilon, and one built on
a set that is not has no reason to come close. Section~\ref{ssec:pdas} measures
the accepted population, whose worst residual is $\qpPdasAcceptWorst$, and reports
what it did and did not observe on the rejecting side. The separation is not
unconditional, and Section~\ref{ssec:transform} closes it deliberately: the
accepted residuals rise with the conditioning, reaching $\qpTfAcceptWorst$ --- four
orders above the figure just quoted --- before candidates begin landing on the far
side of the threshold.
\end{remark}

\subsection{When not to guess}
\label{ssec:whennot}

The certificate makes a rejected guess safe, not free: its cost is paid and
discarded, and the exact walk then runs from the beginning. So there is a
question of policy that the correctness argument does not answer, namely whether
to guess at all, and it turns on an asymmetry. On a problem where the guess would
have been certified, declining to try it forfeits the whole speedup;
on one where it will be rejected, trying costs one guess. At $n = \qpTfTimeN{}$
those two mistakes are worth $\qpTfFalseSkipCost{}$\,ms and
$\qpTfMissedSkipCost{}$\,ms respectively, a factor of $\qpTfCostRatio$
(Section~\ref{ssec:transform}), so a policy should decline only on strong
evidence.

Conditioning is the obvious evidence and turns out to be the wrong kind. It
predicts rejection well in aggregate, but the ratio
$\max_i L_{ii}^2 / \min_i L_{ii}^2$ from the Cholesky factor of $G$ --- free,
since that factorisation is performed regardless --- does not separate the two
populations even ordinally, reading $\qpTfEstGood$ at the worst accepted guess
and $\qpTfEstBad$ at the best rejected one. With the errors that badly ordered
and priced at $\qpTfCostRatio{}$ to one, no threshold on it is worth having.

What is worth having needs no estimate. Across a family of programs sharing $G$
in the sense of Section~\ref{sec:sweep}, a rejection is evidence about the rest
of the family, and it is evidence after the fact rather than a prediction: guess
until the certificate rejects once, and walk thereafter. On the family measured
in Section~\ref{ssec:transform} that recovers $\qpTfFamilySaving$\% of the wasted
work and lands within a percent of the floor set by never guessing at all, while
on a family where guessing works it declines nothing, because the first
rejection never comes. The rule is not a tolerance, so it joins the other two of
this section in not being delicate: it can only cost one guess per family, and it
cannot affect an answer.
